%% P30 --- Cross-entropy and the Kullback--Leibler divergence
%%
%% Stub, generated by tools/gen_stubs.py from tools/programs.json.
%% The brief below is the contract for this program: what it must argue, and
%% what it must buy the reader. Writing the program means deleting the
%% \programstub{} block and replacing it with the program -- after which this
%% file is never regenerated. If the brief turns out to be wrong once the
%% mathematics has been worked, change it in the manifest and record the
%% contradiction in CLAUDE.md under *Resolved questions*.

\program{Cross-entropy and the Kullback--Leibler divergence}
\label{prog:P30}

\programstub{%
Argument: cross-entropy is the cost of coding one distribution with
another's code; KL is the excess; it is non-negative, zero only when the
distributions agree, **and it is not symmetric**. Payoff: the question the
brief asks for \dash{} **what the asymmetry costs you when you pick a
loss.** Forward KL is mode-covering: it pays an unbounded price for putting
no mass where the target has some, so it spreads. Reverse KL is mode-
seeking: it pays for putting mass where the target has none, so it collapses
onto one mode. Distillation minimises one, a variational objective and a KL-
regularised policy the other, and the difference is visible in the output.
Measured on a bimodal target, not asserted. Also: \enquote{KL distance} is a
misnomer and the triangle inequality fails; Jensen--Shannon as the symmetric
alternative and what it costs; and the observation that minimising cross-
entropy against a fixed dataset is minimising forward KL to the empirical
distribution, which closes the loop with P26.

\medskip
\textbf{Planned length:} 55 frames.
}
